% 本文件由 LaTeX 排版 Agent 根据有效 translation.json 自动生成。
% author=Eusebius of Caesarea (c. 265-c. 340); work_id=2804; title=论尼西亚大公会议的信函

\chapter{论尼西亚大公会议的信函}

% structure: p0001 source=h1 target=omitted reason=page-title-already-rendered-by-parent

\Emph{凯撒勒雅的欧瑟伯致其教区人民的信函}。\par

1. 亲爱的诸位，关于在尼西亚召开的大公会议上所议定的教会信仰之事，你们或许已从其他渠道得知，因为传闻往往先于对实际行动的准确描述。但唯恐此类传闻对实情有所歪曲，我们不得不先将我们自己所提出的信仰条文，以及随后教父们在我们的措辞上有所增补而发布的第二条条文，寄送给你们。因此，我们在至虔敬的皇帝面前宣读的那份条文，被宣布为良好且无可指责的，其内容如下：\par

2. \Quote{正如我们从先辈主教们那里所领受的，在我们初次的教理学习中，以及我们领受圣洗时，又正如我们从圣经所学到的，并且我们在长老团中及在主教任上本身所相信和教导的，同样我们如今也如此相信，并在此向你们陈述我们的信仰，即：}\par

3. \Quote{我们相信唯一的天主，全能的圣父，一切有形与无形之物的创造者。并相信唯一的主耶稣基督，天主的圣言，出自天主的天主，出自光明的光明，出自生命的生命，独生子，一切受造物中的首生者，在万世之前由圣父所生，万物藉祂而受造；祂为了我们的救恩而成了血肉，生活在人间，受难，第三日复活，升天，并将再次在光荣中来临，审判生者死者。我们也相信唯一的圣神：}\par

我们相信每一个都是且实存，圣父真实地为圣父，圣子真实地为圣子，圣神真实地为圣神，正如我们的主派遣祂的门徒去宣讲时所说：\Quote{你们去教训万民，因圣父及圣子及圣神之名给他们授洗 \BibleRef{玛 28:19}。}关于祂们，我们坚信我们如此持守，如此思想，且从前已如此持守，并将此信仰持守至死，弃绝一切不虔的异端。我们向全能的天主和我们的主耶稣基督作证，我们从有记忆以来就真心实意地如此思想，如今也真实地如此思想并宣告，而且能够用证据向你们表明并说服你们，即使在过去的时代，我们的信仰和宣讲也是如此。\par

4. 当我们公开提出这一信仰后，似乎没有矛盾之处；但我们的至虔敬皇帝，在任何人之前，首先作证说它包含了最正统的陈述。他还承认这正是他自己的观点，并建议所有在场者同意该信仰，签署其条款并予以赞同，但加上一个词\Quote{\OriginalTerm{同体}{One-in-essence}}。他解释说，这个词并非指身体的属性，也不是说圣子以分裂或割裂的方式从圣父而来；因为那非物质的、理智的、无形体的本性不可能受到任何身体属性的影响，而我们应当以神圣且不可言喻的方式来理解这些事物。这就是我们至智慧、至虔诚的皇帝的神学评述。但教父们为了加入\Quote{同体}一词，起草了以下条文：\par

\Emph{大公会议上所指定的信仰。}\par

\Quote{我们相信唯一的天主，全能的圣父，一切有形与无形之物的创造者：}\par

\Quote{并相信唯一的主耶稣基督，天主的圣子，由圣父所生，独生子，亦即出自圣父的性体；出自天主的天主，出自光明的光明，出自真天主的真天主，受生而非受造，与圣父同体，万物藉祂而受造，无论是在天上的还是在地上的；祂为了我们人类并为了我们的救恩，降下并成了血肉，成了人，受难，第三日复活，升天，并将来审判生者死者。}\par

\Quote{并相信圣神。}\par

\Quote{若有人说\InnerQuote{祂曾不存在}或\InnerQuote{在受生之前祂不存在}或\InnerQuote{祂是从无中而来}，或者有人声称天主子\InnerQuote{出自另一种本体或性体}、或是\InnerQuote{受造的}、\InnerQuote{可变的}、\InnerQuote{可改变的}，天主教会弃绝这样的人。}\par

5. 当他们起草这一条文时，我们并未不加探究地任其通过，而是询问他们引入\Quote{出自圣父的性体}和\Quote{与圣父同体}是何意。因此发生了问题和解释，这些词语的含义经过了理性的审查。他们宣称，\Quote{出自性体}一词表示圣子确实出自圣父，但并非如同圣父的一部分。基于这一理解，我们同意接受这一虔诚教义的含义，因为它教导圣子出自圣父，但并非祂性体的一部分。因此，我们接受了这个含义，甚至没有拒绝\Quote{同体}一词，因为我们将和平以及持守正统观点作为我们的目标。\par

6. 同样，我们也接受了\Quote{受生而非受造}；因为该大公会议声称，\Quote{受造}是其他藉圣子而有的受造物的共同称谓，圣子与它们没有相似之处。因此，他们说，祂不是一件类似于藉祂而来的事物的作品，而是具有一种超越任何作品等级的性体，而神圣的训谕教导说，这性体是由圣父所生，其生育方式对于一切受造的本性而言是深奥难测、不可计算的。\par

7. 同样，经过考察，有理由说圣子与圣父是\Quote{\OriginalTerm{同一实体}{one in essence}}；这并非像有形体那样，也不是像必死之物那样，因为祂不是由本质的分割或分离，也不是由于圣父的本质和德能的任何情感、改变或变化（因为圣父的无始本体与此类一切无关），而是因为\Quote{\OriginalTerm{同一实体}{one in essence}}表明天主子与受造之物毫无相似之处，唯独与生祂的圣父完全相似，并且祂不是出于任何其他的\OriginalTerm{位格}{subsistence}和\OriginalTerm{本质}{essence}，而是出于圣父。这个术语，如此解释之后，我们也认为应当同意；因为我们知道，即使在古代，一些博学而杰出的主教和作家在关于圣父和圣子的神学教导中也曾使用\Quote{\OriginalTerm{同一实体}{one in essence}}这个术语。\par

8. 关于所公布的信德，就说到这里；我们所有人都同意这一信德，并非未经查考，而是根据在至虔诚的皇帝本人面前提出的特定意义，并以上述理由加以证明。至于他们在信德末尾所公布的绝罚条文，并没有使我们感到痛苦，因为它禁止使用圣经中没有的言辞，而教会中几乎所有的混乱和无序都源于这些言辞。既然没有任何天启圣经使用过\Quote{\OriginalTerm{从无中}{out of nothing}}和\Quote{\OriginalTerm{祂曾不存在}{once He was not}}以及随后的其他说法，那么显然没有理由使用或教导这些说法；我们也认为这是一项良好的决定，因为迄今我们并无使用这些说法的习惯。\par

9. 此外，绝罚\Quote{\OriginalTerm{在祂受生之前，祂不存在}{Before His generation He was not}}似乎并不荒谬，因为所有人都承认，天主子在按照肉身的受生之前已经存在。\par

10. 不仅如此，我们至虔诚的皇帝当时还在一次演讲中证明，即使按照祂在万世以前的神圣受生，祂也实际存在，因为甚至在祂有效地受生之前，祂就已经\OriginalTerm{非受生地}{ingenerately}与圣父同在德能中，圣父永远是父亲，如同永远是君王，永远是救主，祂在德能中是一切，并且始终如一、恒常不变。\par

11. 亲爱的诸位，我们被迫将这些传达给你们，以向你们说明我们查考和同意的经过，以及我们如何合理地坚持到最后时刻，只要我们对那些与我们的说法不同的陈述感到反感，但一旦坦率地考察了这些词语的意义，发现它们与我们自己在已公布的信德中所宣认的内容相符，我们就不再痛苦，而是毫无争论地接受了它们。\par

\SourceRecord{Translated by John Henry Newman and Archibald Robertson.；Nicene and Post-Nicene Fathers, Second Series；Vol. 4.；Edited by Philip Schaff and Henry Wace.；Buffalo, NY: Christian Literature Publishing Co.,；1892.；<http://www.newadvent.org/fathers/2804.htm>.}
